\documentclass[10pt,journal]{IEEEtran}

\usepackage{amsmath,amssymb}
\usepackage{cite}
\usepackage{graphicx}
\usepackage{booktabs}
\usepackage{array}
\usepackage{url}

\newcommand{\vect}[1]{\mathbf{#1}}
\newcommand{\mat}[1]{\mathbf{#1}}
\newcommand{\Herm}{\mathrm{H}}
\newcommand{\E}{\mathbb{E}}
\newcommand{\Tr}{\mathrm{Tr}}
\newcommand{\diag}{\mathrm{diag}}
\newcommand{\relu}{\mathrm{ReLU}}

\begin{document}

\title{Communication-Safe Entropy-Guided PPO for Joint Beamforming in Multiuser ISAC Systems}

\author{First~Author,
        Second~Author,
        and~Third~Author%
\thanks{Manuscript received Month Day, 2026; revised Month Day, 2026.}
\thanks{The authors are with the Department of Electrical and Computer
Engineering, University Name, City, Country. Corresponding author:
First Author (e-mail: author@example.com).}
}

\markboth{IEEE Wireless Communications Letters,~Vol.~XX, No.~XX, 2026}%
{Author \MakeLowercase{\textit{et al.}}: Communication-Safe Entropy-Guided PPO for Joint Beamforming}

\maketitle

\begin{abstract}
Joint transmit beamforming is a fundamental design problem in multiuser
integrated sensing and communication (ISAC), where the radar transmit
beampattern must be optimized while the communication users' signal-to-
interference-plus-noise ratio (SINR) constraints are satisfied. Existing
optimization-based methods can provide high-quality solutions, but their
computational cost limits fast inference. This letter proposes a
communication-safe entropy-guided proximal policy optimization (CSE-PPO)
method for ISAC beamforming. The proposed method formulates beamforming as
a finite-horizon residual control problem and introduces a feasibility-
weighted policy entropy that couples exploration with the minimum SINR
margin. Consecutive low-effective-entropy steps are consolidated into
macro-steps for PPO updates, improving credit assignment in low-exploration
regimes. A primal-dual cost update is further used to suppress SINR
violations. Simulation results under the Liu-type multiuser MIMO radar and
communication setup show that CSE-PPO improves radar beampattern quality
over vanilla PPO while maintaining high communication feasibility.
\end{abstract}

\begin{IEEEkeywords}
Integrated sensing and communication, dual-function radar communication,
beamforming, proximal policy optimization, entropy-guided reinforcement
learning.
\end{IEEEkeywords}

\section{Introduction}
\IEEEPARstart{I}{ntegrated} sensing and communication (ISAC) has been
recognized as a key enabler for future wireless systems \cite{liu2022isacsurvey}, where radar
sensing and wireless communication share spectrum, hardware, and spatial
degrees of freedom. A representative multiuser ISAC design is joint
transmit beamforming, in which a colocated transmitter simultaneously
forms radar probing beams and serves multiple downlink users. The classical
formulation minimizes a radar beampattern loss subject to per-antenna
power and user SINR constraints \cite{liu2020joint}. Although semidefinite
relaxation (SDR) can solve this problem with high accuracy, it requires an
optimization problem to be solved for each channel realization.

Learning-based beamforming has recently attracted attention because it can
amortize optimization cost after training \cite{zhao2025learning}. However,
directly applying a neural network or vanilla reinforcement learning (RL)
to constrained ISAC beamforming is challenging. First, unconstrained
exploration can easily violate stringent user SINR constraints. Second,
not all policy updates are equally informative: when the policy entropy is
small, consecutive residual beamforming actions often lead to nearly
indistinguishable rewards and noisy credit assignment.

Motivated by the observation that high-entropy steps drive more effective
RL updates in flow-model alignment \cite{zhang2026egrpo}, this letter
proposes \emph{communication-safe entropy-guided PPO} (CSE-PPO) for
multiuser ISAC beamforming. Different from group-relative policy
optimization, PPO relies on value-based advantage estimates and clipped
policy ratios. Therefore, we adapt the entropy idea through
macro-step PPO updates rather than group-normalized rewards. The main
contributions are summarized as follows:
\begin{itemize}
  \item We formulate joint ISAC beamforming as a finite-horizon residual
  control problem with exact per-antenna power normalization.
  \item We define a feasibility-weighted policy entropy using the minimum
  SINR margin, enabling exploration only when it is communication-safe.
  \item We consolidate consecutive low-effective-entropy steps into
  macro-steps and optimize them with a clipped PPO surrogate.
  \item We provide a reproducible simulation pipeline comparing CSE-PPO
  with SDR, zero-forcing (ZF), vanilla PPO, and entropy-consolidated PPO.
\end{itemize}

\section{System Model and Problem Formulation}
Consider a transmitter equipped with $M$ antennas. It serves $K$ single-
antenna downlink users and simultaneously forms radar beams towards a set
of target directions. Following \cite{liu2020joint}, the transmitted
signal at time $n$ is
\begin{equation}
    \vect{x}[n] = \mat{W}_{c}\vect{c}[n] + \mat{W}_{r}\vect{s}[n],
\end{equation}
where $\mat{W}_{c}\in\mathbb{C}^{M\times K}$ is the communication
beamformer, $\mat{W}_{r}\in\mathbb{C}^{M\times M}$ is the radar
beamformer, $\vect{c}[n]$ denotes the communication symbols, and
$\vect{s}[n]$ denotes the radar waveforms. The transmit covariance matrix
is
\begin{equation}
    \mat{R} = \E[\vect{x}[n]\vect{x}^{\Herm}[n]]
    = \mat{W}_{c}\mat{W}_{c}^{\Herm}
    + \mat{W}_{r}\mat{W}_{r}^{\Herm}.
\end{equation}
The per-antenna power constraint is
\begin{equation}
    [\mat{R}]_{m,m}=P_t/M,\quad m=1,\ldots,M.
\end{equation}

Let $\vect{a}(\theta)$ be the transmit steering vector. The radar
beampattern at direction $\theta$ is
\begin{equation}
    P(\theta;\mat{R})=\vect{a}^{\Herm}(\theta)\mat{R}\vect{a}(\theta).
\end{equation}
Given the desired beampattern $d(\theta)$ and target directions
$\{\theta_p\}_{p=1}^{P}$, the radar loss is
\begin{align}
    L_r(\mat{R},\alpha)
    &= \frac{1}{L}\sum_{\ell=1}^{L}
    \left|\alpha d(\theta_\ell)-P(\theta_\ell;\mat{R})\right|^2  \notag\\
    &\quad + w_c\frac{2}{P^2-P}
    \sum_{p<q}\left|\vect{a}^{\Herm}(\theta_q)
    \mat{R}\vect{a}(\theta_p)\right|^2 .
\end{align}

Let $\mat{H}\in\mathbb{C}^{K\times M}$ be the downlink channel, whose
$k$-th row is denoted by $\vect{h}_k^{\Herm}$. The received SINR of user
$k$ is
\begin{equation}
    \gamma_k =
    \frac{|[\mat{H}\mat{W}_c]_{k,k}|^2}
    {\sum_{i\ne k}|[\mat{H}\mat{W}_c]_{k,i}|^2
    +\sum_{j}|[\mat{H}\mat{W}_r]_{k,j}|^2+\sigma^2 } .
\end{equation}
The joint ISAC beamforming problem is
\begin{subequations}
\begin{align}
    \min_{\mat{W}_c,\mat{W}_r,\alpha}\quad
    & L_r(\mat{R},\alpha)\\
    \mathrm{s.t.}\quad
    & [\mat{R}]_{m,m}=P_t/M,\quad \forall m,\\
    & \gamma_k\ge \Gamma,\quad k=1,\ldots,K,
\end{align}
\end{subequations}
where $\Gamma$ is the required SINR threshold.

\section{Proposed CSE-PPO Method}
\subsection{Residual Beamforming MDP}
We cast the above constrained optimization problem as a finite-horizon
Markov decision process (MDP). Each episode corresponds to one channel
realization and one target-direction configuration. At step $t$, the
state includes the real and imaginary parts of $\mat{H}$, the target
steering vectors, the current beamforming matrix
$\mat{W}_t=[\mat{W}_{c,t},\mat{W}_{r,t}]$, the radar loss, and the user
SINR margins. The policy outputs a real-valued action that is mapped to a
complex residual update $\Delta\mat{W}_t$. The next beamformer is
\begin{equation}
    \widetilde{\mat{W}}_{t+1}=\mat{W}_t+\eta\Delta\mat{W}_t,
\end{equation}
followed by row-wise normalization so that each antenna satisfies the
power constraint exactly:
\begin{equation}
    [\mat{W}_{t+1}]_{m,:}
    =
    \sqrt{\frac{P_t}{M}}
    \frac{[\widetilde{\mat{W}}_{t+1}]_{m,:}}
    {\|[\widetilde{\mat{W}}_{t+1}]_{m,:}\|_2}.
\end{equation}

To reduce harmful exploration, the residual action is parameterized in a
communication-aware form. The communication columns contain diagonal ZF
gain residuals, while both communication and radar columns can use
components from the null space of $\mat{H}$. This structure provides
radar shaping freedom while reducing user interference.

\subsection{Communication-Safe Entropy}
For a Gaussian PPO policy $\pi_\theta(\vect{a}_t|\vect{s}_t)$ \cite{schulman2017ppo}, let
$\mathcal{H}_{\pi,t}$ be the action entropy. Denote the minimum SINR
margin by
\begin{equation}
    m_t=\min_k(\gamma_{k,t}-\Gamma).
\end{equation}
The proposed effective entropy is
\begin{equation}
    \mathcal{H}_{\mathrm{eff},t}
    =
    \frac{\mathcal{H}_{\pi,t}}{d_a}
    \cdot
    \phi(m_t),
    \quad
    \phi(m_t)=
    \frac{1}{1+\exp\{-\beta m_t/\Gamma\}},
\end{equation}
where $d_a$ is the action dimension and $\beta$ controls how sharply the
communication margin affects exploration. When the SINR margin is small,
the effective entropy is reduced even if the policy distribution itself
has large entropy.

\subsection{Entropy-Consolidated PPO Update}
If $\mathcal{H}_{\mathrm{eff},t}$ is above a threshold $\tau$, the step is
updated as a normal PPO transition. Otherwise, consecutive low-effective-
entropy transitions are merged into a macro-step
$\mathcal{B}=\{t,\ldots,t+\ell-1\}$. The macro-step likelihood ratio is
\begin{equation}
    \rho_{\mathcal{B}}(\theta)
    =
    \exp\left(
    \sum_{i\in\mathcal{B}}
    \log\pi_\theta(\vect{a}_i|\vect{s}_i)
    -
    \log\pi_{\theta_{\mathrm{old}}}(\vect{a}_i|\vect{s}_i)
    \right).
\end{equation}
The macro-step advantage is obtained from generalized advantage estimates
(GAE) as
\begin{equation}
    A_{\mathcal{B}}=\sum_{j=0}^{\ell-1}\gamma^j A_{t+j}.
\end{equation}
The clipped PPO objective is then
\begin{equation}
    \mathcal{L}_{\mathrm{clip}}
    =
    \E_{\mathcal{B}}
    \left[
    \min\left(
    \rho_{\mathcal{B}} A_{\mathcal{B}},
    \mathrm{clip}(\rho_{\mathcal{B}},1-\epsilon,1+\epsilon)A_{\mathcal{B}}
    \right)
    \right].
\end{equation}
This design follows the intuition of entropy-aware RL: low-entropy
individual steps often provide weak reward discrimination, while their
merged effect is more informative.

\subsection{Dual-Guided Constraint Control}
The instantaneous reward is
\begin{equation}
    r_t =
    L_r(\mat{W}_{t})-L_r(\mat{W}_{t+1})
    -\xi\|\vect{a}_t\|_2^2 ,
\end{equation}
and the communication violation cost is
\begin{equation}
    c_t=
    \frac{1}{K}\sum_{k=1}^{K}
    \frac{[\Gamma-\gamma_{k,t}]_+}{\Gamma}.
\end{equation}
The PPO advantage is computed using the Lagrangian reward
$\tilde{r}_t=r_t-\lambda c_t$, and the dual variable is updated by
\begin{equation}
    \lambda \leftarrow
    \left[\lambda+\eta_\lambda(\bar{c}-c_0)\right]_+,
\end{equation}
where $\bar{c}$ is the average cost in the rollout batch.

\begin{figure}[t]
\centering
\fbox{\parbox{0.92\linewidth}{
\textbf{CSE-PPO training procedure}\\
1) Initialize beamformer with a ZF-inspired feasible point.\\
2) Collect batched residual beamforming rollouts.\\
3) Compute feasibility-weighted entropy.\\
4) Merge consecutive low-effective-entropy steps.\\
5) Update actor-critic with clipped PPO and dual SINR cost.\\
6) During inference, return the feasible beamformer with the minimum radar loss.
}}
\caption{Communication-safe entropy-guided PPO for ISAC beamforming.}
\label{fig:framework}
\end{figure}

\section{Numerical Results}
We follow the simulation setup in \cite{liu2020joint}. Unless otherwise
specified, the transmitter has $M=10$ antennas, the total transmit power is
$P_t=1$, and the noise variance is $\sigma^2=0.01$. The radar target
directions are $-40^\circ$, $0^\circ$, and $40^\circ$, with an ideal
beamwidth of $10^\circ$. The downlink channel follows Rayleigh fading. The
SINR threshold is varied from $4$ dB to $24$ dB, and the number of users is
selected from $K\in\{2,4,6\}$.

The compared methods are: i) SDR, which solves the relaxed convex problem;
ii) ZF, which uses communication interference suppression and null-space
radar shaping; iii) vanilla PPO; iv) entropy-consolidated PPO (EC-PPO),
which removes the communication feasibility weight; and v) the proposed
CSE-PPO.

\begin{figure}[t]
\centering
\IfFileExists{figs/beampattern.pdf}{%
\includegraphics[width=\linewidth]{figs/beampattern.pdf}
}{%
\fbox{\parbox{0.92\linewidth}{Place generated beampattern figure here. Use
\texttt{plot\_beampatterns.py} to generate \texttt{beampattern.pdf}.}}
}
\caption{Transmit beampatterns obtained by different methods under the same
channel realization.}
\label{fig:beampattern}
\end{figure}

\begin{figure}[t]
\centering
\IfFileExists{figs/tradeoff.pdf}{%
\includegraphics[width=\linewidth]{figs/tradeoff.pdf}
}{%
\fbox{\parbox{0.92\linewidth}{Place radar loss versus SINR threshold figure
here.}}
}
\caption{Radar loss versus communication SINR threshold.}
\label{fig:tradeoff}
\end{figure}

Fig.~\ref{fig:beampattern} shows a representative beampattern comparison.
The SDR method provides the strongest optimization benchmark, while ZF
preserves communication reliability at the cost of radar shaping freedom.
Compared with vanilla PPO, CSE-PPO forms clearer mainlobes around the
target directions and suppresses undesired angular leakage. The gain is
attributed to the communication-safe entropy mechanism, which avoids
wasting policy updates on low-informative and unsafe residual steps.

Fig.~\ref{fig:tradeoff} reports the radar-communication tradeoff. As the
SINR threshold increases, all methods experience radar performance loss.
However, CSE-PPO consistently reduces the gap between learning-based
beamforming and SDR compared with vanilla PPO and EC-PPO. The ablation
between EC-PPO and CSE-PPO validates that weighting entropy by the
communication margin is beneficial for constrained ISAC learning.

\begin{table}[t]
\centering
\caption{Qualitative complexity comparison.}
\label{tab:complexity}
\begin{tabular}{lcc}
\toprule
Method & Online optimization & Inference style\\
\midrule
SDR & High & Convex solver per channel\\
ZF & Low & Closed-form/search\\
Vanilla PPO & Low & Neural forward rollouts\\
EC-PPO & Low & Neural forward rollouts\\
CSE-PPO & Low & Neural forward rollouts\\
\bottomrule
\end{tabular}
\end{table}

Table~\ref{tab:complexity} summarizes the online complexity trend. SDR is
useful as a performance upper benchmark but requires solving a semidefinite
program for each channel. Once trained, PPO-based methods only require a
small number of neural forward passes and deterministic residual updates.

\section{Conclusion}
This letter proposed CSE-PPO for constrained multiuser ISAC beamforming.
By combining residual beamforming, feasibility-weighted entropy,
macro-step PPO updates, and primal-dual SINR control, the proposed method
improves radar beampattern learning while maintaining communication
feasibility. Future work will extend the current framework to wideband
channels and robust beamforming under imperfect CSI.

\bibliographystyle{IEEEtran}
\bibliography{refs}

\end{document}
